Este manual proporciona las instrucciones necesarias para compilar, ejecutar y depurar programas utilizando el sistema \textit{m2r}.

\section{Compilación}
Para compilar el proyecto, asegúrese de tener \texttt{make}, \texttt{gcc}, \texttt{g++}, \texttt{flex} y \texttt{bison} instalados. En la raíz del proyecto, ejecute:
\begin{lstlisting}[language=bash]
make
\end{lstlisting}
Esto generará los ejecutables \texttt{compiler} y \texttt{m2r}.

\section{Ejecución}
El proceso de ejecución se realiza en dos etapas:
\begin{enumerate}
    \item \textbf{Compilación del código fuente:}
    \begin{lstlisting}[language=bash]
./compiler <archivo_fuente.java>
    \end{lstlisting}
    Esto genera un archivo de código intermedio (normalmente \texttt{output.asm}).
    
    \item \textbf{Interpretación:}
    \begin{lstlisting}[language=bash]
./m2r output.asm
    \end{lstlisting}
\end{enumerate}

\section{Depuración en la Máquina Virtual}
El intérprete \texttt{m2r} incluye un depurador básico. Para activarlo (suponiendo que la implementación lo permita mediante banderas o configuraciones), se pueden utilizar comandos interactivos durante la ejecución para:
\begin{itemize}
    \item Listar instrucciones.
    \item Establecer puntos de ruptura (\textit{breakpoints}).
    \item Inspeccionar el contenido de la memoria y registros (Acumulador, Base).
    \item Ejecutar paso a paso (\textit{step-by-step}).
\end{itemize}
\textit{Consulte el código fuente de \texttt{src/m2r.c} para conocer los comandos específicos de depuración soportados.}
